\def\ChapterCode{EXPERT}
\chapter{Antikoagulantien}

Gerinnungshemmende Medikamente sind ein fixer Bestandteil
der Prävention von thromboembolischen Ereignissen bei
Risikopatienten. Dementsprechend stellt sich bei geplanten
und ungeplanten operativen Ereignisen die Frage,w ie mit
diesem Regime umzugehen ist.

Grundsätzlich muss dabei der Grund bzw. das Ziel der
gerinnungshemmenden Therapie sowie das individuelle
Risikoprofil des Patienten einerseits, als auch die Art und
die damit verbundenen potentiellen Blutungskomplikationen
andererseits in Betracht gezogen werden. Weiters zu beachten
ist die Art der Gerinnunghemmung. 


\section{Übersicht der häufigsten gerinnungshemmenden
Medikamente}

\TextSlide
{Coumarine}
{Vitamin-K-Antagonisten

\begin{description}
  \item[Wirkstoffe (Marken)] \T{Phenprocoumon}
  (\I{Marcoumar{\texttrademark}}), \T{Acenocoumarol}
  (\I{Sintrom{\texttrademark}})
  \item[Anwendung] Thromboembolie-Prophylaxe bei Vorhofflimmern,
  Z.\,n. tVT, mechanischer Klappenersatz, biologischer
  Klappenersatz (3 Omnate), Z.\,n. Pulmonalembolie, Z.\,n.
  ACBP
  \item[Laborkontrolle] INR, (Quick)
  \item[Prä-OP-Pause] Jede Unterbrechung erhöht das Risiko
  eines thromboembolischen Ereignisses! In Tabelle
  \ref{tab:bridging} sind mögliche Strategien in
  Abhängigkeit von thromboembolischen und Blutungsrisiko
  dargestellt. Ab einem INR <\,1,5 ist nicht mit einer
  wesentlichen Steigerung des Blutungsrisikos zu rechnen.
  \item[Post-OP-Pause] Siehe \prettyref{tab:bridging}
  \item[Rescue]~
  \begin{itemize}
    \item Semi-akut: Vitamin K (Konakion\texttrademark)
    \Range{2,5}{5}\,--\,10\Mg*
    \item Akut: Fresh Frozen Plasma oder Human prothrombin
    complex (Beriplex{\texttrademark}), Dosierung nach
    Gewicht und INR
  \end{itemize}
\end{description}

}
{
\begin{itemize}
  \item Phenprocoumon
(Marc(o)umar{\texttrademark}) 
\item Acenocoumarol
(Sintrom{\texttrademark})
\end{itemize}
}
{}
{}
{}



\begin{table}[H]
\begin{WidePar}
\caption{Bridging-Therapie von Vitamin-K-Antagonisten mit
LMWH; nach 
\cite{Poldermans2010}
}

\label{tab:bridging}

\begin{tabularx}{\linewidth}{llX}
\addlinespace\toprule\addlinespace
\noindent\TabHeading{Thromboembolisches Risiko}
&
\noindent\TabHeading{Blutungsrisiko}
&
\noindent\TabHeading{Strategie}
\\\addlinespace\midrule\addlinespace
Niedrig 
&
Niedrig 
&
OAK weiter mit INR im therapeutischen Bereich
\\\addlinespace
Niedrig
&
Hoch
&
\begin{compactitem}
  \item OAK 5 Tage prä-OP absetzen
  \item Danach LMWH, letzte Dosis > 12\Hour* prä-OP
  \item 

Post-OP: \begin{compactitem}
  \item LMWH Wiederbeginn ab 12\Hour* bis 2 Tage post-OP, 
  \item OAK Wiederbeginn \Range{1}{2} Tage post-OP, evtl.
  mit Loading Dose
  \item LMWH weiter bis INR im Zielbereich
\end{compactitem}
\end{compactitem}

\Range{1}{2} Tage post-OP, evtl. mit Loading Dose, LMWH
weiter bis INR im Zielbereich
\\\addlinespace
Hoch
&
---
&
\begin{compactitem}
  \item OAK 5 Tage prä-OP absetzen
  \item Danach LMWH in \EE{therapeutischer Dosis}
  (2\,$\times$\,tgl.), letzte Dosis > 12\Hour* prä-OP
  \item 

Post-OP: \begin{compactitem}
  \item LMWH Wiederbeginn ab 12\Hour* bis 2 Tage post-OP, 
  \item OAK Wiederbeginn \Range{1}{2} Tage post-OP, evtl.
  mit Loading Dose
  \item LMWH weiter bis INR im Zielbereich
\end{compactitem}
\end{compactitem}

\Range{1}{2} Tage post-OP, evtl. mit Loading Dose, LMWH
weiter bis INR im Zielbereich
\\
\\\addlinespace\bottomrule\addlinespace

\end{tabularx}

\end{WidePar}
\end{table}



\TextSlide
{Direkter Thrombin Inhibitor (DTI)}
{~

\begin{description}
  \item[Wirkstoffe (Marken)] \T{Dabigatran}
  (\I{Pradaxa{\texttrademark}}), \T{Rivaroxaban}
  (\I{Xarelto{\texttrademark}})
  \item[Anwendung] Orale Thromboseprophylaxe
  nach Hüft- und Kniegelenksersatz,
  Thromboembolieprophylaxe bei Vorhofflimmern;
  Rivaroxaban: Behandlung von tVT, Sekundärprävention von
  tVT und Pulmonalembolien
  \item[Laborkontrolle] eingschränkt; Dabigatran: Ecarin
  Clotting Time; Rivaroxaban: PTZ
  \item[Prä-OP-Pause] 5 Tage, ab Tag -3 LMWH in
  therapeutischer oder prophylaktischer Dosierung bis
  24\Hour* bzw. 12\Hour* vor OP
  \item[Post-OP-Pause] LMWH wieder frühestens
  6\Hour* nach OP, DTI wieder ab Tag 4
%   \item[Rescue]
\end{description}
}
{
\begin{itemize}
  \item Dabigatran (Pradaxa{\texttrademark}), 
  \item Rivaroxaban (Xarelto{\texttrademark}) 
\end{itemize}
}
{}
{}
{}



\TextSlide
{Cox-I-Hemmer}
{
\begin{description}
  \item[Wirkstoffe (Marken)] \T{ASS}
  \item[Anwendung] Chronische stabile KHK, Koronarsklerose,
  Z.\,n. MCI, CABG oder Stentimplantation, Primärprophylaxe
  bei Hochrisikopatienten (Diabetes mellitus, Hypertonie),
  Vorhofflimmern (je nach Risikostratifizierung),
  rheumatische Mitralvitien wenn OAK nicht möglich, u.\,U.
  in Kombination mit OAK
  \item[Laborkontrolle] Eingeschränkt
  \item[Prä-OP-Pause] ASS per se ist keine
  Kontroindikation für chirurgische Eingriffe,
  individuelle Risiko-Nutzen-Analyse. In der
  Praxis \Range{2}{7} Tage. Cave Stent!
  Im nicht-kardiochirurgischen Bereich gibt es keine Evidenz
  bezüglich perioperativen Blutungsrisiko.
  \item[Post-OP-Pause] Uneinheitlich, abhängig von
  internistischer Grunderkrankung
  \item[Rescue] Thrombozytenkonzentrate, Desmopressin
\end{description}
}
{
\begin{itemize}
  \item ASS
\end{itemize}
}
{}
{}
{}



\TextSlide
{ADP-Rezeptor-Antagonisten}
{
\begin{description}
  \item[Wirkstoffe (Marken)] \T{Clopidogrel}
  (\I{Plavix{\texttrademark}}), \T{Ticagrelor}
  (\I{Brilique\texttrademark}), \T{Prasugrel}
  (\I{Efient\texttrademark}), \T{Ticlopidin}
  (\I{Tiklid{\texttrademark}})
  \item[Anwendung] TE-Prophylaxe nach MCI (<\,1 Monat),
  Z.\,n. Insult (<\,6 Monate), pAVK, ACI, Prophylaxe einer
  Stentthrombose nach PCI (Bare Metal Stent (BMS)
  \Range{1}{12} Monate, Drug Eluting Stent (DES)
  \Range{6}{12}Monate)
  \item[Prä-OP-Pause] siehe \ref{topic:Stents}
  und \prettyref{tab:Stents}. Cave Stent! In der Praxis
  \Range{5}{10} Tage (Ticagrelor: 3 Tage). Im
  nicht-kardiochirurgischen Bereich gibt es keine Evidenz bezüglich perioperativen Blutungsrisiko.
  \item[Post-OP-Pause] Je nach Grunderkrankung, evtl. mit
  Loading Dose
  \item[Rescue] Thrombozytenkonzentrate, Desmopressin
\end{description}

}
{
\begin{itemize}
  \item Clopidogrel (Plavix{\texttrademark})
  \item {Ticagrelor}
  ({Brilique\texttrademark})
  \item {Prasugrel}
  ({Efient\texttrademark})
  \item Ticlopidin
  (Tiklid{\texttrademark})
\end{itemize}
}
{}
{}
{}


% \Text
% {}
% {}
% {}
% {}
% {}
% {}


% \begin{table}
% \begin{WidePar}
% \caption{Antidote}
% \small
% \begin{tabularx}{\linewidth}{XXX}
% \addlinespace\toprule\addlinespace
% \noindent\TabHeading{Gruppe}
% &
% \noindent\TabHeading{Wirkstoffe (Marken)}
% &
% \noindent\TabHeading{Häufige Anwendungsgebiete}
% \\\addlinespace\midrule\addlinespace
% \multicolumn{3}{c}{\noindent\TabHeading{Plasmatische
% Gerinnung}} 
% \\\addlinespace\midrule\addlinespace
% Vit.-K-Antagonisten 
% &
% Phenprocoumon (Marc(o)umar{\texttrademark}), Acenocoumarol
% (Sintrom{\texttrademark}) 
% &
% Thromboembolie-Prophylaxe bei Vorhofflimmern, 
% rez. Thrombosen, mechanischer oder biologischer
% Klappenersatz 
% \\\addlinespace
% Direkter Thrombin Inhibitor (DTI)
% &
% Dabigatran (Pradaxa{\texttrademark}), Rivaroxaban
% (Xarelto{\texttrademark}) 
% &
% Alternative zu Vit.-K-Antagonisten
% \\\addlinespace\midrule\addlinespace
% \multicolumn{3}{c}{\noindent\TabHeading{Zelluläre Gerinnung}}
% \\\addlinespace\midrule\addlinespace
% COX-I-Hemmer
% &
% ASS, 
% & Prophylaxe von In-Stent-Thrombosen nach PCI, 
% \\\addlinespace
% ADP-Rezeptor-Antagonisten
% &
% Clopidogrel (Plavix{\texttrademark}), Ticlopidin
% (Tiklid{\texttrademark}), Prasugral (Efient{\texttrademark})
% &
% 
% \\\addlinespace\bottomrule\addlinespace
% 
% \end{tabularx}
% 
% \end{WidePar}
% \end{table}




% 
% \begin{table}
% \begin{WidePar}
% \caption{Antidote}
% 
% \begin{tabularx}{\linewidth}{lXXX}
% \addlinespace\toprule\addlinespace
% \noindent\TabHeading{Agens}
% &
% \noindent\TabHeading{Antidot}
% &
% \noindent\TabHeading{Dosierung}
% &
% \noindent\TabHeading{Anmerkungen}
% \\\addlinespace\midrule\addlinespace
% Vit.-K-Antagonisten 
% &
% Vitamin K (Konnakion{\texttrademark}) 
% &
% 2,5-5\Mg*
% \\\addlinespace
% Vit.-K-Antagonisten
% &
% Fresh Frozen Plasma
% &
% 
% &
% \Range{1000}{1500}\,IE 
% \\\addlinespace
% ASS
% &
% Thrombozytenkonzentrate
% &
% 
% &
% 
% \\\addlinespace
% Heparin
% &
% 
% &
% 
% &
% Wirkungsverlust nach 4\Hour* zu erwarten
% \\\addlinespace
% Heparin
% &
% Protaminsulfat
% &
% 1\Mg* Protaminsulfat pro 100\,U Heparin vor 2\Hour*, max.
% 50\Mg* 
% &
% Für sofortigen Wirkungsverlust; Cave Anaphylaxie!
% \\\addlinespace\bottomrule\addlinespace
% 
% \end{tabularx}
% 
% \end{WidePar}
% \end{table}

\section{Spezielle Fälle}

\subsection{PCI und Stents}

\label{topic:Stents}

Patienten nach einer Koronarintervention stellen ein
besondere Pateintengruppe dar. Nach der Implantation von
Stents besteht ein hohes Risiko einer
\EE{In-Stent-Thrombose}.
Um diesem Phänomen zu begegnen erhalten diese Patienten eine
\EE{duale Plättchenaggregationhemmung} mit \T{ASS} und
\T{Clopidogrel} (\I{Plavix{\texttrademark}}). 

\Cave{Das Absetzen
der dualen Plättchenaggregationshemmung führt zu einer
massiven Erhöhung des Risikos einer In-Stent-Thrombose (90-fach, bis zu 20\PC!). }

Die notwendige
Dauer der dualen Prophylaxe hängt von der Art des Stents ab,
sie beträgt für \I[Bare Metal Stent]{Bare Metal Stents}
(\I{DMS}) mindestens \E{3 Monate}, für \I[Drug Eluting
Stent]{Drug Eluting Stents} (\I{DES}) 12 Monate.
Die Tabelle \ref{tab:Stents} gibt eine Übersicht über die empfohlene
Strategien bei Patienten nach koronarer Intervention, für
Details siehe \cite{PeriopManagementPatMitKoronarstents}. 


\begin{table}[H]
\begin{WidePar}
\caption{Patienten nach PCI; nach ESC/ESA
\citep{Poldermans2010}
}

\label{tab:Stents}

\begin{tabularx}{\linewidth}{lXX}
\addlinespace\toprule\addlinespace
\noindent\TabHeading{Intervention}
&
\noindent\TabHeading{Zeitspanne}
&
\noindent\TabHeading{Strategie}
\\\addlinespace\midrule\addlinespace
\multirow{2}{*}{\noindent\TabSub{Ballondilatation}}
&
< 14 Tage
&
Nicht-dringende OP verschieben
\\\addlinespace
&
$\geq$ 14 Tage
&
OP + ASS
\\\addlinespace\midrule\addlinespace
\multirow{2}{*}{\noindent\TabSub{BMS}}
&
< 6 Wochen
&
Nicht-dringende OP verschieben
\\\addlinespace
&
Min. $\geq$ 6 Wochen

Optimal $\geq$ 3 Monate
&
OP + ASS
\\\addlinespace\midrule\addlinespace
\multirow{2}{*}{\noindent\TabSub{DES}}
&
< 12 Monate
&
Nicht-dringende OP verschieben
\\\addlinespace
&
$\geq$ 12 Monate
&
OP + ASS
\\\addlinespace\bottomrule\addlinespace

\end{tabularx}

\end{WidePar}
\end{table}




\subsection{Lokoregionalanästhesie}

Das bisher Gesagte gilt für Eingriffe in Allgemeinnarkose.
Wird eine Lokoregionalanästhesie (PDA, Plexusblockaden,
\ldots) durchgeführt sind spezielle, z.\,T. strengere Regeln
zu beachten. Details siehe
\cite{RegAnaesthGerinunnungshemmenderMed}. 

\begin{table}[H]
\begin{WidePar}
\caption{Empfohlene Therapiepausen bei
Lokoregionalanästhesie, Auszug; nach ÖGARI
\citep{RegAnaesthGerinunnungshemmenderMed}
}
\label{tab:Lokoregional}
\small

\begin{tabularx}{\linewidth}{lXX}
\addlinespace\toprule\addlinespace
\noindent\TabHeading{Substanz}
&
\noindent\TabHeading{Pause vor Punktion /
Katheteerentfernung} &
\noindent\TabHeading{Beginn nach Punktion /
Katheterentfernung} 
\\\addlinespace\midrule\addlinespace
\noindent\TabSub{LWMH} prophylaktisch
&
11\Hour*
&
2\Hour*
\\
\noindent\TabSub{LWMH} therapeutisch
&
24\Hour* + anti-Xa-Aktivität unter Nachweisgrenze
&
2\Hour*
\\\addlinespace
\noindent\TabSub{Dabigatran}
&
26\Hour*
&
4\Hour*
\\\addlinespace
\noindent\TabSub{Vit.-K-Antagonisten}
&
INR < 1,4 (ca. 2 Tage)
&
sofort
\\\addlinespace
\noindent\TabSub{ASS}
&
nicht erforderlich
&
sofoert
\\\addlinespace
\noindent\TabSub{Clopidogrel}
&
7 Tage
&
sofort
\\\addlinespace
\noindent\TabSub{Ticagrelor}
&
5 Tage
&
6\Hour*
\\\addlinespace
\noindent\TabSub{Prasugrel}
&
\Range{7}{10} Tage
&
6\Hour*
\\\addlinespace\bottomrule\addlinespace

\end{tabularx}

\end{WidePar}
\end{table}

\begin{Literary}
\cite{GerinnungKlinAlltag4}
\fullcite{GerinnungKlinAlltag4}

\cite{Poldermans2010} 
\fullcite{Poldermans2010}

\cite{PeriopManagementPatMitKoronarstents}
\fullcite{PeriopManagementPatMitKoronarstents}

\cite{PeriopBridgingNeueAntikoag}
\fullcite{PeriopBridgingNeueAntikoag}

\cite{RegAnaesthGerinunnungshemmenderMed}
\fullcite{RegAnaesthGerinunnungshemmenderMed}
% 
\end{Literary}

\ChapterEnd